\documentclass[10pt,a4paper]{amsart}
\usepackage[margin=19mm]{geometry}
\usepackage{amsmath,amssymb,mathtools}
\usepackage[scr=rsfs,cal=euler]{mathalfa}
\usepackage{xfrac}
\usepackage[unicode,hidelinks,bookmarks=false]{hyperref}
\newcommand{\ie}{i.e.\ }
\newcommand{\cf}{cf.\ }
\newcommand{\resp}{respectively\ }
\newcommand{\Subsection}{\subsection}
\providecommand{\nobreakdash}{\nobreak\hbox{-}}
\setlength{\emergencystretch}{2em}
\multlinegap0pt
\usepackage{amssymb}
\usepackage{mathtools}
\usepackage[shortlabels]{enumitem}
\newtheorem{theorem}{Theorem}
\newtheorem{lemma}{Lemma}
\usepackage{cleveref}
\crefname{theorem}{Theorem}{Theorems}
\crefname{lemma}{Lemma}{Lemmas}
\newcommand{\Pw}{\overline{\mathcal P}}
\newcommand{\ZZ}{\mathbb Z}
\title{Real polynomials with given multiplicities of real roots: Complete conjectural description of homology}
\author{Boris Shapiro}
\date{Source: arXiv:2608.19733v2 (CC BY 4.0)}
\begin{document}
\maketitle

\noindent\emph{Source and licence.} Real polynomials with given multiplicities of real roots: Complete conjectural description of homology. Version arXiv:2608.19733v2 (CC BY 4.0), distributed by arXiv under the Creative Commons Attribution 4.0 International licence, http://creativecommons.org/licenses/by/4.0/, as recorded in the arXiv licence metadata for this exact version and shown on the abstract page https://arxiv.org/abs/2608.19733v2. Full licence notice: \texttt{licenses/arxiv-2608.19733-CC-BY-4.0.txt}.\par\smallskip

\noindent\textit{Editorial scope.} Counted unit: the article's theorem thm:second-main (source0.tex lines 122--133). The article's complete original proof of it is delivered with it (source0.tex lines 592--619, one \texttt{proof} environment of 28 lines, which the article places away from the statement). No mathematics is written, repaired or re-derived: the statement and the proof are the article's own lines, in the article's own notation, and no retained line is substituted or re-flowed.  Retained uncounted as the source-local context the proof needs: lines 203--217; lines 558--561; lines 572--582.  Nothing was omitted from any retained span, so the package carries no omission record.  No retained line contains a citation, so the wrapper carries no bibliography and no bibliography entry.  No retained line contains a figure environment, an image-inclusion command or drawing code, so no image is delivered and the package declares no asset dependency.  Editorial formatting only, in addition: an AMS \texttt{amsart} preamble that restates the article's own declarations and macros used by the retained lines, namely the environments \texttt{theorem}, \texttt{lemma} on separate counters, together with the standard packages those lines need.  The retained environments print in this document as Theorem 1, Lemma 1 in order of appearance, whereas the article prints the same items with the numbering of its own full document.  The complete native source of the article as distributed in the arXiv e-print for this exact version is bundled unchanged as \texttt{sources/208120/source0.tex}.
\begin{lemma}[Membership criterion]\label{lem:membership}\label{lm:2.1}
Let $\omega=(\omega_1,\dots,\omega_s)$ and
$\eta=(\eta_1,\dots,\eta_r)$.  Then $\eta\preceq\omega$ if and only if there
are indices
\[
  0=i_0\le i_1\le\cdots\le i_r=s
\]
such that, for every $j=1,\dots,r$,
\[
  \eta_j\ge \sum_{k=i_{j-1}+1}^{i_j}\omega_k,
  \qquad
  \eta_j\equiv \sum_{k=i_{j-1}+1}^{i_j}\omega_k\pmod2.
\]
The sum is understood as $0$ when $i_{j-1}=i_j$.
\end{lemma}

\begin{equation}\label{eq:3115-euler-series}
 \frac{F_{(3,1,1,5)}(t)}{1-t^2}
 =-\frac{u^9(1-u)}{(1-u^2)(1-u^3)}.
\end{equation}

\begin{lemma}[Integral unit cancellation]\label{lem:unit-cancellation}
Let $C$ be a finite based free chain complex.  Suppose $u\in C_k$ and
$v\in C_{k-1}$ are basis vectors and the coefficient
$a=[v]\partial u$ is $1$ or $-1$.  Deleting $u,v$ and replacing the
differential on the remaining generators by
\[
 D'_{y,x}=D_{y,x}-D_{y,u}a^{-1}D_{v,x}
 \qquad(x,y\notin\{u,v\})
\]
produces an integrally chain-homotopy-equivalent complex.
\end{lemma}

\begin{theorem}[Cancellation counterexample]\label{thm:second-main}
For $\omega=(3,1,1,5)$ and $d=24$,
\[
  \widetilde H_k\bigl(\Pw_{24}^{\langle3,1,1,5\rangle};\ZZ\bigr)
  \cong
  \begin{cases}
    \ZZ,&k=7,8,\\
    0,&k\notin\{7,8\}.
  \end{cases}
\]
Consequently $\chi_{24}(3,1,1,5)=0$ although the reduced homology is nonzero.
\end{theorem}

\begin{proof}[Computer-assisted proof of \Cref{thm:second-main}]
The ideal $\langle3,1,1,5\rangle_{24}$ was generated independently in two
ways: by direct closure under merges and insertions, and by the finite-state
membership criterion of \Cref{lem:membership}.  Both enumerations give exactly
$76\,384$ cells.  Their numbers by length are
\[
\begin{array}{c|rrrrrr}
r&1&2&3&4&5&6\\\hline
\#&8&100&655&2817&8174&15832
\end{array}
\qquad
\begin{array}{c|rrrrr}
r&7&8&9&10&11\\\hline
\#&20265&16857&8760&2586&330.
\end{array}
\]
The alternating sum is zero, independently confirming
\eqref{eq:3115-euler-series} in degree $24$.

Every repeated incidence in the boundary was aggregated before reduction and
$\partial^2=0$ was checked on every basis vector.  A replay of $38\,191$
integral unit cancellations of the form in
\Cref{lem:unit-cancellation} leaves exactly two isolated generators, one in
degree $7$ and one in degree $8$, with zero residual differential.  Hence the
original complex is integrally chain-homotopy equivalent to
$\ZZ[7]\oplus\ZZ[8]$, which proves the theorem and rules out all torsion and
all additional homology.
\end{proof}

\end{document}
